% GENERATED FILE. Edit canonical lessons, then run npm run generate:books.
% canonical-source-hash: fnv1a64:48c6f0475bd0a60f
% canonical-lessons: TE-C213-akasmattuga, TE-C213-pade-pade, TE-C213-kasepu, TE-C213-sekanu, TE-C213-monna

\chapter{Suddenly, Again and again, Little while, Second, Day before yesterday}
\label{ch:te-suddenly-again-and-again-little-while-second-day-before-yesterday}
\hlchaptermodality{213}

\begin{chapteropening}
\textbf{By the end of this chapter:} \emph{I can say suddenly, again and again, a little while, a second and the day before yesterday.}

Complete the last lesson of chapter 213: I can say suddenly, again and again, a little while, a second and the day before yesterday.
\end{chapteropening}

\section[\texorpdfstring{\te{అకస్మాత్తుగా} (akasmāttugā)}{అకస్మాత్తుగా (akasmāttugā)}]{\te{అకస్మాత్తుగా} (akasmāttugā) — suddenly}

\label{lesson:TE-C213-akasmattuga}



\begin{quote}
Before the new one: say the Telugu for until, then the Telugu for Sankranti.
\end{quote}

\subsection*{You'll want to know: \te{అకస్మాత్తుగా}}
\textbf{\te{అకస్మాత్తుగా}} — \emph{akasmāttugā} — \textquotedblleft{}suddenly\textquotedblright{}.

\begin{grammarlens}[title={What you've built}]
One more word for this chapter.
\end{grammarlens}

\subsection*{Your turn}
\begin{itemize}
\raggedright
  \item \emph{Say it:} \emph{akasmāttugā}
  \item \emph{Say it:} \emph{akasmāttugā}, once more
  \item \emph{Say it:} say \emph{akasmāttugā}
  \item \emph{Recall:} say the Telugu for until, then the Telugu for Sankranti, then say \emph{akasmāttugā} again
\end{itemize}

\begin{tcolorbox}[breakable,colback=teal!4,colframe=teal!35!black,title={Before you move on}]
What does \te{అకస్మాత్తుగా} mean? (\textquotedblleft{}Suddenly\textquotedblright{}.) Say it once more.
\end{tcolorbox}
\section[\texorpdfstring{\te{పదే} \te{పదే} (padē padē)}{పదే పదే (padē padē)}]{\te{పదే} \te{పదే} (padē padē) — again and again}

\label{lesson:TE-C213-pade-pade}



\begin{quote}
Before the new one: say the Telugu for Sankranti, then the Telugu for suddenly.
\end{quote}

\subsection*{You'll want to know: \te{పదే} \te{పదే}}
\textbf{\te{పదే} \te{పదే}} — \emph{padē padē} — \textquotedblleft{}again and again\textquotedblright{}.

\begin{grammarlens}[title={What you've built}]
One more word for this chapter.
\end{grammarlens}

\subsection*{Your turn}
\begin{itemize}
\raggedright
  \item \emph{Say it:} \emph{padē padē}
  \item \emph{Say it:} \emph{padē padē}, once more
  \item \emph{Say it:} say \emph{padē padē}
  \item \emph{Recall:} say the Telugu for Sankranti, then the Telugu for suddenly, then say \emph{padē padē} again
\end{itemize}

\begin{tcolorbox}[breakable,colback=teal!4,colframe=teal!35!black,title={Before you move on}]
What does \te{పదే} \te{పదే} mean? (\textquotedblleft{}Again and again\textquotedblright{}.) Say it once more.
\end{tcolorbox}
\section[\texorpdfstring{\te{కాసేపు} (kāsēpu)}{కాసేపు (kāsēpu)}]{\te{కాసేపు} (kāsēpu) — a little while}

\label{lesson:TE-C213-kasepu}



\begin{quote}
Before the new one: say the Telugu for suddenly, then the Telugu for again and again.
\end{quote}

\subsection*{You'll want to know: \te{కాసేపు}}
\textbf{\te{కాసేపు}} — \emph{kāsēpu} — \textquotedblleft{}a little while\textquotedblright{}.

\begin{grammarlens}[title={What you've built}]
One more word for this chapter.
\end{grammarlens}

\subsection*{Your turn}
\begin{itemize}
\raggedright
  \item \emph{Say it:} \emph{kāsēpu}
  \item \emph{Say it:} \emph{kāsēpu}, once more
  \item \emph{Say it:} say \emph{kāsēpu}
  \item \emph{Recall:} say the Telugu for suddenly, then the Telugu for again and again, then say \emph{kāsēpu} again
\end{itemize}

\begin{tcolorbox}[breakable,colback=teal!4,colframe=teal!35!black,title={Before you move on}]
What does \te{కాసేపు} mean? (\textquotedblleft{}A little while\textquotedblright{}.) Say it once more.
\end{tcolorbox}
\section[\texorpdfstring{\te{సెకను} (sekanu)}{సెకను (sekanu)}]{\te{సెకను} (sekanu) — a second}

\label{lesson:TE-C213-sekanu}



\begin{quote}
Before the new one: say the Telugu for again and again, then the Telugu for a little while.
\end{quote}

\subsection*{You'll want to know: \te{సెకను}}
\textbf{\te{సెకను}} — \emph{sekanu} — \textquotedblleft{}a second\textquotedblright{}.

\begin{grammarlens}[title={What you've built}]
One more word for this chapter.
\end{grammarlens}

\subsection*{Your turn}
\begin{itemize}
\raggedright
  \item \emph{Say it:} \emph{sekanu}
  \item \emph{Say it:} \emph{sekanu}, once more
  \item \emph{Say it:} say \emph{sekanu}
  \item \emph{Recall:} say the Telugu for again and again, then the Telugu for a little while, then say \emph{sekanu} again
\end{itemize}

\begin{tcolorbox}[breakable,colback=teal!4,colframe=teal!35!black,title={Before you move on}]
What does \te{సెకను} mean? (\textquotedblleft{}A second\textquotedblright{}.) Say it once more.
\end{tcolorbox}
\section[\texorpdfstring{\te{మొన్న} (monna)}{మొన్న (monna)}]{\te{మొన్న} (monna) — the day before yesterday}

\label{lesson:TE-C213-monna}



\begin{quote}
Before the new one: say the Telugu for a little while, then the Telugu for a second.
\end{quote}

\subsection*{You'll want to know: \te{మొన్న}}
\textbf{\te{మొన్న}} — \emph{monna} — \textquotedblleft{}the day before yesterday\textquotedblright{}.

\begin{grammarlens}[title={What you've built}]
One more word for this chapter.
\end{grammarlens}

\subsection*{Your turn}
\begin{itemize}
\raggedright
  \item \emph{Say it:} \emph{monna}
  \item \emph{Say it:} \emph{monna}, once more
  \item \emph{Say it:} say \emph{monna}
  \item \emph{Recall:} say the Telugu for a little while, then the Telugu for a second, then say \emph{monna} again
\end{itemize}

\begin{tcolorbox}[breakable,colback=teal!4,colframe=teal!35!black,title={Before you move on}]
What does \te{మొన్న} mean? (\textquotedblleft{}The day before yesterday\textquotedblright{}.) Say it once more.
\end{tcolorbox}
